% ============================================================
%  investor_content.tex --- Descripcion General (Inversionista)
%  Proyecto: Conjunto Residencial Cerrado (65 casas)
%  Reemplazar Curiti Postmoderno con el nombre real
% ============================================================

% ─────────────────────────────────────────────────────────────
\section{Descripci\'on General del Proyecto}
% ─────────────────────────────────────────────────────────────

\begin{darkbox}
  {\large\bfseries\color{GoldAccent}
    Curiti Postmoderno: Infraestructura El\'ectrica de Nueva
    Generaci\'on para Vivienda Sostenible\par}
  \vspace{6pt}
  El proyecto consiste en el dise\~no e implementaci\'on de la
  infraestructura el\'ectrica completa de un conjunto residencial
  cerrado de 65 viviendas unifamiliares en Bucaramanga, Colombia.
  Su diferencial estrat\'egico radica en tres pilares t\'ecnicos:
  distribuci\'on subterr\'anea de alta confiabilidad, carga nativa
  para veh\'iculos el\'ectricos en cada unidad, y un sistema de
  macromedici\'on en media tensi\'on que elimina las p\'erdidas
  no controladas de energ\'ia en las zonas comunes.
\end{darkbox}

\vspace{10pt}

% ─────────────────────────────────────────────────────────────
\section{Propuesta de Valor T\'ecnica}
% ─────────────────────────────────────────────────────────────

\subsection{Infraestructura El\'ectrica Subterr\'anea}

La totalidad de la red de distribuci\'on del conjunto opera en
media tensi\'on a $13{,}2\text{ kV}$, con conductores tipo THHN
instalados en canaletas subterr\'aneas. Esta decisi\'on de dise\~no
elimina la exposici\'on a condiciones ambientales adversas, reduce
la frecuencia de mantenimientos correctivos y mejora
significativamente la est\'etica urbana del proyecto,
factores que inciden directamente en la valoraci\'on de los
inmuebles.

La topolog\'ia adoptada es \textbf{radial con selectividad por
caja de protecci\'on}, lo que significa que una falla el\'ectrica
en cualquier vivienda o sector queda aislada localmente sin
afectar al resto del conjunto. Esta caracter\'istica,
frecuentemente ausente en proyectos de menor especificaci\'on,
reduce los tiempos de interrupci\'on del servicio y los costos
operativos de mantenimiento para la administraci\'on.

\begin{insightbox}[Ventaja frente a la competencia]
La red subterr\'anea con topolog\'ia radial selectiva representa
un est\'andar superior al m\'inimo exigido por el Reglamento
T\'ecnico de Instalaciones El\'ectricas (RETIE). Proyectos
comparables en el segmento medio-alto de Bucaramanga operan con
redes a\'ereas convencionales, lo que coloca a \textbf{Curiti Postmoderno} en una categor\'ia diferenciada de calidad
percibida y real.
\end{insightbox}

\subsection{Infraestructura Nativa para Veh\'iculos El\'ectricos}

Cada una de las 65 viviendas cuenta con \textbf{dos acometidas
trif\'asicas independientes} y dos medidores individuales
instalados desde la etapa constructiva:

\begin{center}
\begin{tabular}{
  >{\bfseries}l
  >{\centering\arraybackslash}p{4cm}
  p{6.5cm}}
  \toprule
  Medidor & Uso & Descripci\'on \\
  \midrule
  Medidor 1 &
    Consumo residencial &
    Registro del consumo dom\'estico
    convencional de la vivienda \\
  \rowcolor{NavyDeep!6!white}
  Medidor 2 &
    Cargador de VE &
    Registro exclusivo del consumo del
    cargador de veh\'iculo el\'ectrico,
    con cuenta independiente ante ESSA \\
  \bottomrule
\end{tabular}
\end{center}

\vspace{6pt}

La separaci\'on de cuentas entre consumo residencial y carga
de veh\'iculo el\'ectrico permite al propietario acceder a
\textbf{tarifas diferenciadas} que las comercializadoras aplican
para usuarios con alto consumo, y facilita la declaraci\'on de
gastos de movilidad en contextos empresariales. Para el
inversionista, este atributo posiciona el proyecto en el mercado
de compradores con veh\'iculo el\'ectrico o en proceso de
adopci\'on, un segmento de alto poder adquisitivo en crecimiento
sostenido en Colombia.

% ─────────────────────────────────────────────────────────────
\section{Sistema de Medici\'on y Eficiencia Energ\'etica}
% ─────────────────────────────────────────────────────────────

\subsection{Macromedici\'on en Media Tensi\'on}

El esquema de medici\'on implementado constituye el elemento
t\'ecnico m\'as innovador del proyecto desde la perspectiva
financiera. Se instal\'o un \textbf{macromedidor general en
la frontera del proyecto a $13{,}2\text{ kV}$}, que registra
la totalidad de la energ\'ia que ingresa al conjunto.
Aguas abajo, los 130 medidores individuales (2 por vivienda)
registran el consumo de cada unidad privada.

El consumo de las zonas comunes --- alumbrado, bombas
hidr\'aulicas, administraci\'on --- se obtiene mediante un
\textbf{balance matem\'atico de energ\'ia}:

\vspace{4pt}
\begin{center}
\begin{tabular}{c}
  \toprule
  \large$E_{\text{zonas comunes}} =
    E_{\text{frontera MT}} -
    \displaystyle\sum E_{\text{medidores BT}}$ \\[6pt]
  \bottomrule
\end{tabular}
\end{center}
\vspace{4pt}

Este mecanismo, conocido t\'ecnicamente como
\textbf{balance de energ\'ia por macromedici\'on}, elimina
la posibilidad de que existan p\'erdidas de energ\'ia
no contabilizadas (\'energia ``negra'') que en proyectos
convencionales terminan siendo asumidas por la
administraci\'on --- y en \'ultima instancia por los
propietarios --- sin ninguna visibilidad ni control.

\begin{insightbox}[Impacto financiero para la administraci\'on]
En conjuntos residenciales sin macromedici\'on, las p\'erdidas
no contabilizadas en redes de distribuci\'on interna suelen
representar entre el 3\% y el 8\% del consumo total facturado,
costo que se traslada a la cuota de administraci\'on. El sistema
de balance de energ\'ia de \textbf{Curiti Postmoderno}
hace que cada kWh sea trazable, reduciendo este rubro a cero
y fortaleciendo la transparencia financiera ante los copropietarios.
\end{insightbox}

\subsection{Legalidad y Tr\'amite ante ESSA}

El esquema de macromedici\'on con balance de energ\'ia ha sido
dise\~nado para su \textbf{radicaci\'on formal ante la
Electrificadora de Santander (ESSA)}, la comercializadora
de energ\'ia de la regi\'on. La legalidad del esquema est\'a
respaldada por la normativa colombiana vigente y el proceso
de legalizaci\'on se ejecuta en la etapa de entrega del
proyecto, garantizando que la facturaci\'on de zonas comunes
quede configurada por diferencia de lecturas desde el primer
d\'ia de operaci\'on.

Esta gesti\'on proactiva ante la comercializadora protege
al conjunto residencial de contingencias regulatorias futuras
y le otorga al inversionista la certeza de que la
infraestructura entregada es completamente legal, auditable
y sostenible en el largo plazo.

% ─────────────────────────────────────────────────────────────
\section{Cumplimiento Normativo}
% ─────────────────────────────────────────────────────────────

El dise\~no el\'ectrico del proyecto fue desarrollado con
estricto apego al marco regulatorio colombiano vigente:

\begin{tabular}{ll}
  \toprule
  \bfseries Norma & \bfseries Alcance \\
  \midrule
  RETIE (Res. 90708 de 2013 y actualizaciones) &
    Dise\~no, materiales y seguridad \\
  NTC 2050 &
    C\'odigo El\'ectrico Nacional \\
  Reglamento de Distribuci\'on ESSA &
    Conexi\'on y medici\'on \\
  Ley 1715 de 2014 &
    Energ\'ias renovables y eficiencia \\
  \bottomrule
\end{tabular}

\vspace{8pt}

El proceso de revisi\'on cruzada entre los miembros del equipo
de dise\~no garantiz\'o la consistencia entre memorias de
c\'alculo, plantillas de verificaci\'on y planos constructivos
antes de la entrega formal, reduciendo el riesgo de
observaciones por parte de los organismos de inspecci\'on.

% ─────────────────────────────────────────────────────────────
\section{Resumen de Atributos para el Inversionista}
% ─────────────────────────────────────────────────────────────

\begin{tabularx}{\linewidth}{
  >{\bfseries}X
  >{\centering\arraybackslash}p{2.8cm}
  X}
  \toprule
  Atributo & Impacto & Beneficio concreto \\
  \midrule
  Red 100\% subterr\'anea THHN &
    Alto &
    Mayor vida \'util, cero mantenimiento a\'ereo,
    est\'etica valorizada \\
  Topolog\'ia radial selectiva &
    Medio-alto &
    Fallas aisladas; no se apaga el conjunto
    por un evento puntual \\
  2 medidores por vivienda &
    Alto &
    Cuenta VE independiente; acceso a tarifas
    diferenciadas para el propietario \\
  Macromedici\'on MT + balance &
    Alto &
    Cero p\'erdidas negras; cuota de administraci\'on
    transparente y controlada \\
  Legalidad ante ESSA &
    Cr\'itico &
    Sin contingencias regulatorias;
    facturaci\'on operativa desde el d\'ia 1 \\
  Cumplimiento RETIE pleno &
    Cr\'itico &
    Habilitaci\'on sin observaciones;
    reducci\'on de riesgo en inspecci\'on \\
  \bottomrule
\end{tabularx}

\vspace{12pt}
\begin{darkbox}
  {\bfseries\color{GoldAccent}
    Una infraestructura que agrega valor real\par}
  \vspace{6pt}
  \textbf{Curiti Postmoderno} no solo cumple los requisitos
  m\'inimos de la norma colombiana: los supera. La combinaci\'on
  de red subterr\'anea, carga nativa para veh\'iculos el\'ectricos
  y un sistema de medici\'on transparente en la frontera del
  proyecto representa una propuesta de valor diferenciada,
  auditable y sostenible en el tiempo --- los tres atributos
  que un inversionista exige antes de comprometer capital.
\end{darkbox}